% ================================================================
\section{Summary: The Physics--Mathematics Dictionary}
\label{sec:summary}
% ================================================================

The following table summarizes the complete physics--mathematics
dictionary of the Cathedral, mapping every formal object in the
Lean~4 codebase to its physical counterpart. The full phenomenon
map (120+ entries, organized by physical domain) appears in
Appendix~A (Section~\ref{app:phenomenon-map}).

\begin{table}[h]
\centering
\small
\begin{tabular}{@{}>{}p{4cm}>{}p{5cm}>{}p{4.5cm}@{}}
\toprule
\textbf{Cathedral Object} & \textbf{Physics Concept} & \textbf{Lean Module} \\
\midrule
$L^2(0,1)$ & Hilbert space (Fock space) & \lean{Defs.lean} \\
$\mathbf{1} \in L^2$ & Vacuum state $\ket{0}$ & --- \\
$h_k(x) = \fract{1/(kx)}$ & Single-particle state & \lean{Defs.lean} \\
$f_{N,v}(x) = \sum v_k h_k$ & Trial wavefunction & \lean{BDMellin.lean} \\
$d_N^2 = \inf_v \|1 - f_v\|^2$ & Ground state energy & \lean{NymanBeurling.lean} \\
$G(j,k) = \braket{h_j}{h_k}$ & Two-point correlator & \lean{Gram/*.lean} \\
$G > 0$ & Reflection positivity & \lean{Gram/Bounds.lean} \\
$C = G - bb^T$ & Covariance / fluctuations & \lean{Vasyunin/*.lean} \\
$\lambda_{\min}(G)$ & Mass gap / spectral gap & \lean{Spectral/*.lean} \\
$\lambda_{\min} \sim c/\ln N$ & Mass gap closing & \lean{Sieve/ParitySchur.lean} \\
Mod-8 decomposition & Internal symmetry sectors & \lean{Spectral/ClassRestriction.lean} \\
Schur complement & Supersymmetric localization & \lean{Structural/SchurComplement.lean} \\
\midrule
\multicolumn{3}{@{}l@{}}{\emph{Vasyunin Tower (Lattice Field Theory, \S\ref{sec:vasyunin})}} \\
\midrule
Per-tile FTC & Per-site Lagrangian & \lean{CrossTermFTC.lean} \textbf{(proved)} \\
Row partition & Kadanoff block-spin grouping & \lean{OffDiagPartition.lean} \textbf{(proved)} \\
Telescoping sums & Transfer matrix propagation & \lean{TelescopeSum.lean} \textbf{(proved)} \\
Stirling approx & Saddle-point / steepest descent & \lean{StirlingBridge.lean} \textbf{(proved)} \\
Gauss digamma formula & Bloch decomposition ($\ZZ/q\ZZ$) & \lean{DigammaReflection.lean} \textbf{(axiom)} \\
Convergence axiom & Ergodic hypothesis & \lean{ConvergenceAxioms.lean} \textbf{(axiom)} \\
Beatty sequence bound & Nearest-neighbor interaction & \lean{CrossTermFTC.lean} \textbf{(proved)} \\
$V(a,b) = \sum\fract{mb/a}\cot(\pi m/a)$ & Scattering phase shift & \lean{FormulaBridge.lean} \\
Dedekind reciprocity & Time-reversal invariance & \lean{LogDigammaBridge.lean} \textbf{(axiom)} \\
Uniqueness of limits & Ergodic = thermodynamic & \lean{LogDigammaBridge.lean} \textbf{(proved)} \\
\midrule
\multicolumn{3}{@{}l@{}}{\emph{Parseval Bridge and Mertens}} \\
\midrule
Parseval identity & Unitarity / optical theorem & \lean{PlancherelBypass.lean} \\
Three-term decomposition & Tree + loop + self-energy & \lean{PlancherelBypass.lean} \\
$1/|s|^2$ integral $= 1$ & Cauchy/Breit-Wigner resonance & \lean{PlancherelBypass.lean} \\
$M(x) = \sum \mu(n)$ & Magnetization & \lean{MertensBound.lean} \\
$\mu(n) \in \{-1,0,1\}$ & Spin variable & --- \\
RH $\Rightarrow |M(x)| = O(\sqrt{x})$ & Critical exponent $\beta = 1/2$ & \lean{MertensBound.lean} \\
Log-cutoff taper & Bartlett window / UV cutoff & \lean{MertensIntegral.lean} \\
Abel $S_1, S_2, S_3$ & Magnetization / $\chi$ / $c_v$ & \lean{AbelTail/*.lean} \textbf{(proved)} \\
$1 - b^Tv$ decomposition & Thermodynamic hierarchy & \lean{CalcBounds.lean} \textbf{(proved)} \\
\midrule
\multicolumn{3}{@{}l@{}}{\emph{Perron Chain (Inverse Laplace, \S\ref{sec:perron})}} \\
\midrule
Perron contour integral & Inverse Laplace transform & \lean{Perron/KernelBound.lean} \textbf{(proved)} \\
Perron half-integer assembly & Quantitative inverse Laplace & \lean{Perron/HalfIntegerPerron.lean} \textbf{(proved)} \\
Born--Oppenheimer $\|M-A_N\|+\|A_N-B\|$ & Fast/slow mode separation & \lean{Perron/HalfIntegerPerron.lean} \textbf{(proved)} \\
Archimedean $N$ choice & Adaptive UV regularization & \lean{Perron/HalfIntegerPerron.lean} \textbf{(proved)} \\
$h_{\mathrm{tail\_crushed}}$ & Asymptotic freedom & \lean{Perron/HalfIntegerPerron.lean} \textbf{(proved)} \\
Contour shift $\sigma \to 1/2$ & Crossing phase boundary & \lean{Perron/Rectangle.lean} \textbf{(proved)} \\
\midrule
\multicolumn{3}{@{}l@{}}{\emph{Analytic Tools}} \\
\midrule
Borel--Carath\'eodory & Maximum entropy principle & \lean{(Mathlib)} \\
Hadamard Three-Circles & Open-closed / conformal gauge & \lean{Zeta/Hadamard.lean} \textbf{(proved)} \\
Zero-counting axiom & Weyl law / spectral density & \lean{Zeta/Hadamard.lean} \textbf{(axiom)} \\
$|\zeta(s)| \geq c/|t|^A$ & Free energy lower bound & \lean{Zeta/Convexity.lean} \textbf{(proved)} \\
RH $\Rightarrow \zeta(s) \neq 0$ & No off-shell resonances & \lean{Zeta/Convexity.lean} \textbf{(proved)} \\
$C \approx 21.649$ & $1/c_v$ (inverse heat capacity) & \lean{(Rust oracle, 512-bit)} \\
$c_{\text{holes}} \approx 0.046$ & Casimir energy / sum rule & --- \\
$\theta(t)$ functional equation & Modular invariance & \lean{ThetaBound.lean} \\
$\Lambda(s) < 0$ on $(0,1)$ & Vacuum instability at real $s$ & \lean{BDMellin.lean} \\
$x = e^{-u}$ substitution & Wick rotation & \lean{White/Kinematics.lean} \\
Weyl's inequality & Eigenvalue perturbation theory & \lean{RayleighBridge.lean} \\
Sherman--Morrison & Rank-1 scattering update & \lean{ShermanMorrison.lean} \\
PNT axioms $S_1, S_2, S_3$ & Response function asymptotics & \lean{PNT/AbelMean.lean} \\
\bottomrule
\end{tabular}
\caption{The complete physics--mathematics dictionary of the Cathedral
(updated May~10, 2026, v17 Oracle Capstone). The table is organized by subsystem.
Entries marked \textbf{(proved)} denote theorems with zero sorry;
\textbf{(axiom)} denotes off-crown mathematical axioms.
\textbf{One} literature axiom remains on the primary crown path;
alternative paths achieve the same result with two to four axioms;
the Oracle Bridge achieves RH with one trusted computation axiom.
The forward chain from RH to $d_N^2 \to 0$ is now a continuous,
compiler-verified logical chain with zero sorry on the crown path.
Every mathematical equation in the proof corresponds to a physical
concept, and this correspondence is not metaphorical but structural.}
\label{tab:dictionary}
\end{table}

